% !TEX root = ../algorithms.tex

\section{Вероятностные алгоритмы и хэш-таблицы}

\subsection{Вероятностные алгоритмы}

Вероятностный алгоритм~--- это алгоритм, использующий генератор случайных чисел
\texttt{random}.

\begin{Definition}{Функция random}{random}
	Функция $\mathrm{random}(i,j)\colon \mathbb{R}\times\mathbb{R}\to\mathbb{R}$
	возвращает равновероятно выбранное значение $x\in\{i,\,i{+}1,\ldots,j\}$.
\end{Definition}

\begin{Remark}{}{}
	На практике в языках программирования используют \emph{псевдогенераторы}.
	В данном курсе будем считать, что генератор истинно случаен.
	Современные криптографические инструкции (аппаратные генераторы шума) приближаются
	к истинной случайности.
\end{Remark}

\begin{Example}{Поиск позиции нуля в массиве}{find-zero}
	\textbf{Вход:} массив $a[1,\ldots,n]$, в котором не менее половины элементов равны~$0$.
	\textbf{Задача:} найти позицию~$0$.
	
	\begin{verbatim}
		do {
			p = random(1, n);
		} while (a[p] != 0);
		return p;
	\end{verbatim}
	
	На каждой итерации вероятность успеха не менее $\tfrac{1}{2}$,
	поэтому ожидаемое число итераций не превышает~$2$.
\end{Example}

Пусть $t(a)$~--- время работы алгоритма на конкретном входе~$a$.
Величина $t(a)$ случайна (зависит от выборов \texttt{random}).

\begin{Definition}{Ожидаемое время работы}{expected-time}
	\emph{Ожидаемое время работы} алгоритма на входах длины~$n$:
	\[
	T(n) \;=\; \max_{a\,\in\, I_n} \mathbf{E}\bigl[t(a)\bigr],
	\]
	где $I_n$~--- множество всех входов длины~$n$.
\end{Definition}

Для алгоритма поиска нуля:
\[
T(n) \;\leq\; O\!\left(\sum_{z=1}^{\infty} z\cdot\frac{1}{2^z}\right) = O(1) \;\ll\; O(n).
\]

\subsubsection*{Классы вероятностных алгоритмов}

\begin{Definition}{Алгоритм Лас-Вегас}{las-vegas}
	\emph{Алгоритм Лас-Вегас}~--- вероятностный алгоритм, который \textbf{всегда}
	возвращает корректный ответ, но время работы~$T(n)$ является случайной величиной.
\end{Definition}

\begin{Definition}{Алгоритм Монте-Карло}{monte-carlo}
	\emph{Алгоритм Монте-Карло}~--- вероятностный алгоритм, который может вернуть
	\textbf{некорректное} решение, однако вероятность ошибки мала и поддаётся оценке.
\end{Definition}

\begin{Remark}{}{}
	Алгоритмы Лас-Вегас встречаются чаще.
	Алгоритмы Монте-Карло важны, когда детерминированное решение слишком дорого.
\end{Remark}

\subsection{Алгоритм Фрейвальдса: проверка произведения матриц}

\begin{Example}{Проверка $AB = C$}{freivalds}
	\textbf{Вход:} $A, B, C \in \mathbb{R}^{n\times n}$.
	Верно ли, что $AB = C$?
	
	Детерминированные алгоритмы: $O(n^{\omega})$, $\omega > 2{,}37$.
	
	\textbf{Алгоритм Фрейвальдса} (Монте-Карло, $O(n^2)$):
	\begin{enumerate}
		\item Выбрать $\delta\in\{0,1\}^n$ равновероятно и покоординатно независимо.
		\item Вычислить $AB\delta$ и $C\delta$.
		\item Если $AB\delta = C\delta$, вернуть \textbf{<<да>>}; иначе~--- \textbf{<<нет>>}.
	\end{enumerate}
\end{Example}

\begin{Theorem}{Оценка ошибки алгоритма Фрейвальдса}{freivalds-err}
	Если $AB = C$, алгоритм всегда отвечает верно.
	Если $AB \neq C$, то $\Pr(\text{ошибка}) \leq \tfrac{1}{2}$.
\end{Theorem}

\begin{proof}
	Случай $AB = C$ очевиден. Пусть $AB \neq C$.
	Пусть матрицы отличаются в строке~$k$; обозначим $x := (AB)[k,\cdot]$,
	$y := C[k,\cdot]$, тогда $x \neq y$.
	Ошибка~--- событие $\{x\delta = y\delta\}$, т.\,е.\ $(x-y)\delta = 0$.
	
	Пусть $\ell$~--- индекс, в котором $x_\ell \neq y_\ell$.
	Зафиксируем все $\delta_i$ при $i \neq \ell$.
	Условие $(x-y)\delta = 0$ записывается как
	\[
	(x_\ell - y_\ell)\,\delta_\ell \;=\; -\!\sum_{i \neq \ell}(x_i - y_i)\,\delta_i.
	\]
	\begin{enumerate}
		\item Если правая часть равна~$0$: равенство выполняется только при $\delta_\ell = 0$,
		поэтому $\Pr(x\delta = y\delta \mid \delta_{-\ell}) = \tfrac{1}{2}$.
		\item Если правая часть $\neq 0$: уравнение имеет не более одного решения
		$\delta_\ell \in \{0,1\}$, поэтому $\Pr(x\delta = y\delta \mid \delta_{-\ell}) \leq \tfrac{1}{2}$.
	\end{enumerate}
	В обоих случаях $\Pr(\text{ошибка}) \leq \tfrac{1}{2}$.
\end{proof}

\subsubsection*{Уменьшение вероятности ошибки}

	Чтобы добиться $\Pr(\text{ошибка}) < \varepsilon$, достаточно запустить алгоритм
	$k$ раз независимо, где $k \in \mathbb{N}$ удовлетворяет $\dfrac{1}{2^k} < \varepsilon$.
	Если хотя бы в одном запуске ответ \textbf{<<нет>>}~--- вернуть \textbf{<<нет>>};
	если во всех~--- \textbf{<<да>>}~--- вернуть \textbf{<<да>>}.
	Тогда $\Pr(\text{ошибка}) \leq \dfrac{1}{2^k} < \varepsilon$.

\begin{Remark}{}{}
	Исторически метод Монте-Карло применялся для вычисления интегралов и решения
	дифференциальных уравнений: случайные точки $(p_x, p_y)$ накидываются равномерно,
	и доля точек с $p_y \leq f(p_x)$ приближает $\int f(x)\,dx$.
\end{Remark}

\subsection{Проверка тождественного равенства многочленов}

\begin{Definition}{Задача PIT}{pit}
	Даны многочлены $p(x_1,\ldots,x_n)$ и $q(x_1,\ldots,x_n)$ над полем~$\mathbb{F}$.
	Проверить: $p \equiv q$, то есть $\forall\,x_1,\ldots,x_n\colon p = q$.
	Эквивалентно: $f \equiv 0$, где $f := p - q$.
\end{Definition}

	Если $f$ задан вычислительной схемой размера $O(k)$, то $\deg f \leq 2^k$.
	При $|\mathbb{F}| = +\infty$ детерминированный алгоритм мог бы раскрыть~$f$
	в сумму мономов
	$f(x_1,\ldots,x_n) = \sum_{\mathbf{i}} a_{\mathbf{i}}\,x_1^{d_1}\cdots x_n^{d_n}$
	и проверить все $a_{\mathbf{i}} = 0$, однако это экспоненциально от размера входа.
	В детерминированном случае никто лучше не может.
	
	Для Монте-Карло алгоритма: выбрать случайное $a \in \mathbb{F}^n$
	и проверить $f(a) \neq 0$.

\subsection{Лемма Шварца--Зиппеля}

\begin{Definition}{Степень многочлена от нескольких переменных}{poly-degree}
	\[
	\deg f \;=\; \max\Bigl\{\,\textstyle\sum_{j} d_j \;\Big|\; x_1^{d_1}\cdots x_n^{d_n}
	\text{ --- моном } f\Bigr\}.
	\]
\end{Definition}

\begin{Lemma}{Шварца--Зиппеля}{sz}
	Пусть $f(x_1,\ldots,x_n) \not\equiv 0$ над полем~$\mathbb{F}$,
	$S \subseteq \mathbb{F}$, $0 < |S| < +\infty$.
	Если $x_1,\ldots,x_n$ выбраны равновероятно и независимо из~$S$, то
	\[
	\Pr_{x_1,\ldots,x_n \in S}\!\bigl(f(x_1,\ldots,x_n) = 0\bigr) \;\leq\; \frac{\deg f}{|S|}.
	\]
\end{Lemma}

\begin{proof}
	Индукция по $n$.
	
	\textbf{База} $n = 1$. Многочлен $f(x_1)$ имеет не более $\deg f$ корней, поэтому
	$\Pr_{x_1 \in S}(f(x_1) = 0) \leq \deg f / |S|$.
	
	\textbf{Шаг} $n > 1$. Запишем $f(x_1,\ldots,x_n) = \sum_{k=0}^{\deg f} p_k(x_1,\ldots,x_{n-1})\cdot x_n^k$.
	Так как $f \not\equiv 0$, существует наибольшее $k'$ такое, что $p_{k'} \not\equiv 0$.
	По индукционному предположению:
	\[
	\Pr\bigl(p_{k'}(x_1,\ldots,x_{n-1}) \neq 0\bigr) \;\geq\; 1 - \frac{\deg p_{k'}}{|S|}.
	\]
	Зафиксируем $x_1,\ldots,x_{n-1}$ так, что $p_{k'} \neq 0$.
	Тогда $g(x_n) = \sum_k p_k(\ldots)\,x_n^k$ является ненулевым многочленом от~$x_n$
	степени $\leq \deg f - \deg p_{k'}$. По базе:
	\[
	\Pr\bigl(g(x_n) \neq 0\bigr) \;\geq\; 1 - \frac{\deg f - \deg p_{k'}}{|S|}.
	\]
	Объединяя оба условия:
	\begin{align*}
		\Pr\bigl(f \neq 0\bigr)
		&\;\geq\; \Bigl(1 - \frac{\deg f}{|S|}\Bigr)\Bigl(1 - \frac{\deg f - \deg p_{k'}}{|S|}\Bigr) \\
		&= 1 - \frac{\deg f}{|S|} - \frac{\deg f - \deg p_{k'}}{|S|}
		+ \frac{\deg p_{k'} \cdot (\deg f - \deg p_{k'})}{|S|^2}
		\;\geq\; 1 - \frac{\deg f}{|S|}.
	\end{align*}
	Следовательно $\Pr(f = 0) \leq \deg f / |S|$.
\end{proof}

\begin{Remark}{}{}
	С помощью леммы Шварца--Зиппеля задача PIT решается Монте-Карло алгоритмом:
	выбрать точку из $S^n$ и проверить $f = 0$.
	Удобно считать поле бесконечным~--- тогда вероятность ошибки стремится к нулю
	с ростом~$|S|$.
\end{Remark}

\subsection{Хэш-таблицы}

\subsubsection*{Словари: определение и операции}

\begin{Definition}{Словарь}{dict}
	\emph{Словарь}~--- структура данных для хранения пар «ключ--значение»
	($x \in U$, значение $\delta$), поддерживающая операции:
	\begin{itemize}
		\item $\mathrm{insert}(x,\,\delta)$~--- вставить значение $\delta$ по ключу~$x$;
		\item $\mathrm{delete}(x)$~--- удалить запись с ключом~$x$;
		\item $\delta = \mathrm{query}(x)$~--- найти значение по ключу~$x$.
	\end{itemize}
	Здесь $U$~--- \emph{универсальное множество} всех возможных ключей (большое, например $\mathbb{Z}$).
\end{Definition}

	Все три операции выполняются за амортизированное $O(1)$.

\subsubsection*{Хэширование произвольных ключей}

Если $U$~--- не целые числа, возможны два подхода.

\begin{enumerate}
	\item \textbf{Уникальные объекты}
	($x = y \Leftrightarrow$ это один и тот же объект в памяти).
	Каждому объекту сопоставляется адрес в памяти $\in \mathbb{N}$
	(в C++ это буквально указатель).
	
	\item \textbf{Неуникальные объекты} (строки, кортежи, точки, \ldots).
	Сопоставляем им \emph{хэш} (например, хэш Карпа--Рабина).
	Коллизия~--- событие редкое, с малой вероятностью.
\end{enumerate}